O comércio de moda pela internet oferece catálogos com dezenas de milhares de peças, e quem compra raramente consegue ver todas. Sistemas de recomendação reduzem essa escolha a uma lista curta, ordenada pela chance de cada item interessar à pessoa~\cite{deshpande2004}. Em moda, porém, uma blusa é escolhida para ser usada com uma calça, um sapato e uma bolsa. A pergunta passa a ser o que combina com aquilo que a pessoa já tem~\cite{han2017,vasileva2018}. Este trabalho trata dessa versão do problema. Dado um look incompleto e a categoria da peça que falta, o sistema ordena todas as peças dessa categoria, e a peça certa deve ficar entre as primeiras.

O problema aparece em situações concretas. Na página de uma saia, uma loja virtual pode sugerir o sapato que completa o conjunto. Um aplicativo que organiza o armário pode indicar qual peça fecha o look do dia. Nos dois casos, o sistema precisa responder rápido e justificar a sugestão, por exemplo com looks montados por outras pessoas em que as peças aparecem juntas. A tarefa é difícil: na base usada, cada consulta tem em média 10.940 candidatas, e o sorteio coloca a peça certa entre as 10 primeiras em só 0,182\% dos casos.

A equipe escolheu o tema pelo interesse em aplicações de moda e pela chance de aplicar as estruturas da disciplina a dados reais. Com uma entrada desse tamanho, a diferença entre um algoritmo eficiente e um ingênuo fica visível.

Os trabalhos de referência tratam a compatibilidade entre peças com redes neurais treinadas sobre imagens~\cite{han2017,vasileva2018,mcauley2015}, e as notas dessas redes são difíceis de explicar. Optamos por estruturas clássicas por três motivos: elas rodam em C++ com a biblioteca padrão, têm custo conhecido e permitem explicar cada nota pelos looks reais, pelas peças parecidas e pelos grupos de estilo.

O objetivo principal é construir e avaliar um recomendador que completa looks sobre a base FashionVerse-MME~\cite{fashionverse}, com 102.370 peças e 31.333 looks. O recomendador usa KD-tree, grafo de peças e looks com busca em largura (BFS), Union-Find, Kruskal e Tarjan. Os objetivos específicos são: (i) verificar se a combinação da semelhança entre peças com os looks reais acerta mais que cada fonte isolada e mais que a recomendação pelas peças mais populares; (ii) medir como o tempo de cada estrutura cresce de $10^3$ a $10^5$ peças e comparar esse crescimento com a teoria.

As contribuições são as seguintes. Primeiro, uma sequência de versões, da A até a versão final F. A F soma a nota da versão C, os pontos entre grupos formados pelo Kruskal e a popularidade. Segundo, uma avaliação justa: todas as peças da categoria recebem nota, cada resultado é a média de 10 ordens de desempate, os pesos são escolhidos só na validação e as diferenças vêm com intervalo de 95\% por bootstrap~\cite{efron1979}. Terceiro, a base tratada, publicada no Zenodo~\cite{zenodo}. Por fim, o código aberto, com instruções para reproduzir os resultados~\cite{repositorio}.

A Seção~II apresenta o estado da arte e as estruturas usadas, com a complexidade de cada uma. A Seção~III descreve os dados, a modelagem, o algoritmo e a configuração dos experimentos. A Seção~IV compara as versões, separa a diferença real da que pode ser sorte e verifica se os tempos medidos seguem a teoria. A Seção~V resume contribuições, limitações e próximas etapas. Esperamos mostrar que estruturas clássicas acertam bem acima do acaso quando a peça procurada já apareceu em looks de treino, e deixar claro onde elas falham.
